\begin{proposition}[3.1.6]
\label{IV.2.3.1.6-fr}
Sean $X$ un preesquema localmente noetheriano y $\mathcal F$ un
$\mathcal O_X$-módulo coherente; entonces $\operatorname{Ass}(\mathcal F)$ es
localmente finito \textup{(es decir, todo punto de $X$ admite una vecindad cuya
intersección con $\operatorname{Ass}(\mathcal F)$ es finita)}.
\end{proposition}

Basta considerar el caso en que $X$ es afín, y por tanto noetheriano; entonces
la proposición resulta de (3.1.2) y de Bourbaki, \emph{Alg. comm.}, cap.~IV,
\S~1, núm.~4, cor. del teo.~2.

\oldpage[IV-2]{38}
\begin{proposition}[3.1.7]
\EGAReaderTarget{ega4.c04.prop.3-1-7}{IV.2.3.1.7-fr}
Sea $X$ un preesquema.
\begin{enumerate}
\item[\emph{(i)}] Sea $0\to\mathcal F'\to\mathcal F\to\mathcal F''\to0$
una sucesión exacta de $\mathcal O_X$-módulos cuasicoherentes. Entonces
$\operatorname{Ass}(\mathcal F')\subset\operatorname{Ass}(\mathcal F)
\subset\operatorname{Ass}(\mathcal F')\cup\operatorname{Ass}(\mathcal F'')$.
\item[\emph{(ii)}] Sean $\mathcal F$ un $\mathcal O_X$-módulo
cuasicoherente y $(\mathcal F_\alpha)$ una familia de
sub-$\mathcal O_X$-módulos cuasicoherentes de $\mathcal F$ tal que
$\mathcal F$ sea la unión de los $\mathcal F_\alpha$. Entonces
$\operatorname{Ass}(\mathcal F)=\bigcup_\alpha\operatorname{Ass}(\mathcal F_\alpha)$.
\item[\emph{(iii)}] Para toda familia $(\mathcal F_\alpha)$ de
$\mathcal O_X$-módulos cuasicoherentes, se tiene
$\operatorname{Ass}(\bigoplus_\alpha\mathcal F_\alpha)
=\bigcup_\alpha\operatorname{Ass}(\mathcal F_\alpha)$.
\end{enumerate}
\end{proposition}

Se reduce de inmediato a las proposiciones correspondientes para los módulos
(Bourbaki, \emph{loc. cit.}, \S~1, núm.~1, fórmula (1), prop.~3 y cor.~1 de la
prop.~3).

\begin{proposition}[3.1.8]
\label{IV.2.3.1.8-fr}
Sean $X$ un preesquema localmente noetheriano, $\mathcal F$ un
$\mathcal O_X$-módulo cuasicoherente, $U$ un abierto de $X$ y $\mathcal J$ un
ideal coherente de $\mathcal O_X$ que define un subpreesquema cerrado de $X$
cuyo espacio subyacente es $X\setmin U$. Las condiciones siguientes son equivalentes:
\begin{enumerate}
\item[\emph{a)}] $\operatorname{Ass}(\mathcal F)\subset U$.
\item[\emph{b)}] Para todo abierto afín $V$, toda sección de $\mathcal F$
sobre $V$ cuya restricción a $V\cap U$ sea nula es igual a $0$.
\item[\emph{c)}] El homomorfismo canónico
$\mathcal F\to\mathcal Hom_{\mathcal O_X}(\mathcal J,\mathcal F)$ es
inyectivo.
\end{enumerate}
\end{proposition}

Como la cuestión es local, se puede suponer $X=\operatorname{Spec}(A)$ afín,
siendo $A$ un anillo noetheriano, $\mathcal F=\widetilde M$, donde $M$ es un
$A$-módulo, y $\mathcal J=\widetilde{\mathfrak J}$, donde $\mathfrak J$ es un
ideal de $A$. El homomorfismo $M\to\operatorname{Hom}_A(\mathfrak J,M)$ hace
corresponder a $m\in M$ el homomorfismo $x\leadsto xm$ de $\mathfrak J$ en
$M$; decir que no es inyectivo significa que existe un $m\ne0$ en $M$ tal que
$\mathfrak Jm=0$.

Demostremos primero que \emph{c)} implica \emph{a)}: en efecto, si
$\operatorname{Ass}(\mathcal F)$ intersecara entonces $X\setmin U$, existiría un
ideal primo $\mathfrak p\in\operatorname{Ass}(M)$ que contiene a
$\mathfrak J$, y por tanto un elemento $m\ne0$ de $M$ tal que
$\mathfrak Jm=0$. En segundo lugar, \emph{b)} implica \emph{c)}: en efecto,
si se tuviera $\mathfrak Jm=0$ para un $m\ne0$ en $M$, para todo ideal primo
$\mathfrak q\not\supset\mathfrak J$ existiría un $a\in\mathfrak J$ que no
pertenece a $\mathfrak q$; por tanto, la relación $am=0$ implicaría que la
imagen canónica de $m$ en $M_\mathfrak q$ es $0$; en otras palabras, $m$
sería una sección $\ne0$ de $\mathcal F$ sobre $X$ cuya restricción a $U$
sería nula.

Finalmente, \emph{a)} implica \emph{b)}. Para verlo, observemos que el
homomorfismo canónico
$M\longrightarrow\prod_{\mathfrak p\in\operatorname{Ass}(M)}M_\mathfrak p$
es \emph{inyectivo}: en efecto, si $N$ es el núcleo de este homomorfismo, se
tiene $\operatorname{Ass}(N)\subset\operatorname{Ass}(M)$; si existiera un
$\mathfrak p\in\operatorname{Ass}(N)$, habría un elemento $n\in N$ cuyo
anulador sería $\mathfrak p$; pero, por definición de $N$, existiría un
elemento $s\notin\mathfrak p$ tal que $sn=0$, lo cual es absurdo; se concluye
que $\operatorname{Ass}(N)=\varnothing$, y por tanto $N=0$. Ahora bien, la
condición \emph{a)} implica que, si $m\in M$ es una sección de $\mathcal F$
cuya restricción a $U$ es nula, la imagen canónica de $m$ en
$M_\mathfrak p$ es nula para todo $\mathfrak p\in\operatorname{Ass}(M)$, y
por tanto $m=0$. Q.E.D.

